%% 致谢

\begin{acknowledgements}
毕业设计与论文写作即将结束。回顾本课题从选题确定、理论推导、实验实现到论文定稿的整个过程，我得到了许多老师、同学、朋友和家人的帮助与支持。在此，我向所有给予我指导、帮助和鼓励的人表示诚挚的感谢。

首先，衷心感谢我的校内指导老师以及上海人工智能实验室的各位老师和同学。本次毕业设计的主要工作在实验室完成，期间不仅使用了实验室提供的计算资源，也受益于实验室围绕大语言模型后训练和强化学习对齐方向的日常讨论。各位老师在方法构思、理论分析和实验设计等方面给予了我重要指导，帮助我发现并修正了研究过程中存在的问题；实验室的同学们也在代码实现、训练调试和实验结果整理等方面给予了耐心帮助。一起排查实验异常、分析训练曲线和讨论实验结果的经历，使我对大语言模型偏好对齐研究有了更加具体和深入的理解。

感谢学院各位老师在本科四年中的培养与教导。系统的专业课程学习为我打下了计算机科学与技术方面的基础，也使我能够在毕业设计中综合运用算法、机器学习、深度学习和系统实现等相关知识。

感谢身边的同学和朋友。在论文写作和实验推进过程中，我曾多次遇到瓶颈，大家的讨论、鼓励和陪伴帮助我逐渐理清思路，并保持稳定的研究节奏。也感谢家人长期以来的理解与支持，使我能够专注于学业和研究。

最后，感谢评阅老师和答辩委员会各位老师对本文的审阅与指导。由于本人能力和时间有限，本文仍存在不足之处，恳请各位老师批评指正。我也将在今后的学习和研究中继续改进，不断提升自己的专业能力和科研水平。
\end{acknowledgements}
